\chapter{Embasamento Teórico}
\section{Habilidades Sociais}
De acordo com \citeautor{caballo2003manual}, as HS envolvem componentes comportamentais, cognitivos e fisiológicos, abrangendo tanto a expressão verbal e não verbal quanto os processos de percepção e interpretação social. O déficit em HS está associado a \textbf{dificuldades de comunicação}, \textbf{isolamento social} e maior vulnerabilidade a \textbf{transtornos emocionais}, enquanto um bom repertório de HS contribui para \textbf{relações interpessoais saudáveis}, \textbf{bem-estar psicológico} e \textbf{sucesso profissional}.

O THS se baseia em técnicas como modelagem, ensaios de comportamento, \textit{feedback} e reforço positivo, visando desenvolver competências como iniciar e manter conversas, defender direitos pessoais, lidar com críticas e expressar emoções de maneira assertiva \cite{caballo2003manual}. Essas estratégias tornam o THS uma ferramenta versátil, utilizada em contextos clínicos, escolares e organizacionais.

No contexto contemporâneo, a interação social também se manifesta nos ambientes virtuais, como nos jogos \textit{multiplayer online}, que favorecem a cooperação, a comunicação e a resolução conjunta de problemas. Pesquisas recentes têm investigado o potencial desses jogos como espaços de treinamento informal de habilidades sociais, especialmente entre jovens \cite{granic2014}. Essa perspectiva sugere uma possível convergência entre os princípios do THS e a dinâmica dos ambientes virtuais, hipótese que será explorada ao longo deste trabalho.
%\todo{tirei o TDAH}
\iffalse
\section{TDAH}
O \textit{Transtorno de Déficit de Atenção e Hiperatividade} (TDAH) é classificado como um transtorno do neurodesenvolvimento, caracterizado por um padrão persistente de \textbf{desatenção}, \textbf{hiperatividade} e \textbf{impulsividade}. Esses sintomas devem ser inapropriados para o nível de desenvolvimento do indivíduo e causar prejuízos significativos em múltiplos contextos (por exemplo, em casa, na escola ou no trabalho) \cite{dsm-5}.

Pesquisas indicam que o TDAH possui um forte componente genético com alterações neurobiológicas, especialmente em regiões como o \textit{córtex pré-frontal}, responsável pelo controle executivo, atenção e inibição de respostas \cite{barkley1994attention}. Além disso, fatores ambientais, como estresse familiar e adversidades na infância, podem atuar como moduladores do quadro.

Indivíduos com TDAH frequentemente apresentam \textbf{dificuldades sociais}, como problemas na interpretação de sinais sociais, maior propensão a conflitos interpessoais e baixa tolerância à frustração. A combinação desses fatores pode gerar \textbf{prejuízos acadêmicos, sociais e emocionais} ao longo da vida, reforçando a importância de intervenções precoces e integradas \cite{barkley1994attention}.

Segundo \citeautor{barkley1994attention} o tratamento do TDAH é geralmente \textbf{multimodal}, envolvendo:
\begin{itemize}
    \item Intervenções farmacológicas;
    \item Terapias psicossociais e comportamentais;
    \item Estratégias educacionais.
\end{itemize}

\fi
\section{Jogos Multiplayer}

Os \textbf{jogos digitais} evoluíram de simples experimentos tecnológicos para uma indústria de grande relevância cultural e social. Além do entretenimento, a literatura aponta benefícios cognitivos, emocionais e sociais proporcionados por essas experiências interativas.  \citeautor{granic2014} destacam que os jogos digitais influenciam positivamente quatro domínios principais: (i) \textit{funções cognitivas}, como memória de trabalho, atenção seletiva e pensamento estratégico; (ii) \textit{motivação intrínseca}, que é reforçada pela estrutura de desafios e recompensas; (iii) \textit{regulação emocional}, com jogos sendo usados para lidar com estresse e ansiedade; e (iv) \textit{habilidades sociais}, promovendo interação, cooperação e comunicação.

Já os jogos \textit{multiplayer}  têm ganhado destaque no meio acadêmico como ambientes interativos capazes de promover habilidades sociais em contextos diversos. A natureza cooperativa ou competitiva desses jogos exige que os jogadores se comuniquem, colaborem e lidem com conflitos em tempo real, oferecendo oportunidades práticas de exercício de competências como empatia, escuta ativa e trabalho em equipe.

Diversos estudos têm explorado esse potencial em diferentes faixas etárias e contextos. Pesquisas recentes mostram que jogos como Minecraft e plataformas de realidade aumentada ou virtual, por exemplo, promovem interações sociais significativas e facilitam o desenvolvimento de comportamentos pró-sociais \cite{Jagannath2020, Bailey2021}.

Além do contexto terapêutico e educacional, ambientes \textit{multiplayer} também demonstram potencial em promover apoio emocional e senso de pertencimento, como observado em comunidades de jogos online \cite{Martoncik2016}. Esses espaços virtuais funcionam muitas vezes como zonas seguras de socialização para indivíduos que enfrentam dificuldades em interações presenciais, proporcionando experiências sociais que contribuem para o bem-estar emocional.

Dessa forma, os jogos multiplayer se apresentam como ferramentas promissoras para intervenções voltadas ao desenvolvimento de habilidades sociais, por combinarem engajamento lúdico com interações sociais significativas. Essa combinação justifica sua escolha neste trabalho como meio experimental para investigar possíveis benefícios sociais em jogadores.